% Generated draft: requires mathematical and source review.
\documentclass[11pt]{article}
\usepackage[margin=25mm]{geometry}
\usepackage{fontspec}
\setmainfont{Latin Modern Roman}
\usepackage{amsmath,amssymb,mathtools,unicode-math}
\setmathfont{Latin Modern Math}
\usepackage{enumitem,needspace,xurl,hyperref}
\setlength{\parindent}{0pt}
\setlength{\parskip}{5pt}
\setlength{\emergencystretch}{3em}
\newcommand{\N}{\mathbb{N}}
\newcommand{\Z}{\mathbb{Z}}
\newcommand{\Q}{\mathbb{Q}}
\newcommand{\R}{\mathbb{R}}
\newcommand{\C}{\mathbb{C}}
\newcommand{\F}{\mathbb{F}}
\newcommand{\A}{\mathbb{A}}
\newcommand{\PP}{\mathbb{P}}
\newcommand{\E}{\mathbb{E}}
\newcommand{\eps}{\varepsilon}
\newcommand{\dd}{\,\mathrm{d}}
\begin{document}
\section*{UnsolvedMath ID: 30003978}
\textbf{Source problem code:} catalogue entry
\subsection{Definitions and mathematical statement}
Let \(r\geq 9\) be an integer, and let
\[
\pi:X_r\longrightarrow \mathbb P^2_k
\]
be the blowup at \(r\) distinct general points. Write \(H\) for the pullback of the class of a line and \(E_1,\ldots,E_r\) for the exceptional divisors. The supplied record does not specify the ground field \(k\), its characteristic, or the precise meaning of generality of the centers; these data require confirmation.

For an ample line bundle \(L\) on \(X_r\) and a geometric point \(x\in X_r\), define the Seshadri constant by
\[
\varepsilon(L;x)=\inf_{\substack{C\subset X_r\ \text{an integral curve}\\x\in C}}
\frac{c_1(L)\cdot C}{\operatorname{mult}_x C}.
\]
Here \(\operatorname{mult}_x C\) is the multiplicity of \(C\) at \(x\). The phrase ``at a very general point'' means that the asserted property holds outside a countable union of proper Zariski closed subsets, in a ground-field setting where this interpretation is meaningful.

One standard non-strict numerical formulation of the Nagata hypothesis is
\[
\operatorname{Nag}_{\mathrm{nef}}(X_r):\qquad
\sqrt r\,H-\sum_{i=1}^r E_i\ \text{is nef}.
\]
A real divisor class is nef if its intersection with every integral curve is nonnegative. Equivalently, this hypothesis requires
\[
\sqrt r\,\deg B\ \geq\ \sum_{i=1}^r\operatorname{mult}_{p_i}B
\]
for every integral plane curve \(B\), where \(p_1,\ldots,p_r\) are the blowup centers.

\textbf{Question, conditional on this interpretation of the Nagata hypothesis.} For every integer \(r\geq 9\), does
\[
\operatorname{Nag}_{\mathrm{nef}}(X_r)
\quad\Longrightarrow\quad
\exists\,L\in\operatorname{Pic}(X_r)\ \text{ample such that}
\varepsilon(L;x)\in\mathbb R\setminus\mathbb Q
\ \text{for very general }x\in X_r
\]
hold, assuming only the Nagata conjecture for this same \(X_r\)? The line bundle is chosen before the very general point. No Nagata hypothesis for a different number of blowup centers, or additional conjectural hypothesis, is supplied.

The source does not state the precise version of the Nagata conjecture it intends. Consequently, the displayed implication is a formulation template requiring confirmation of that hypothesis and the ground-field conventions. In particular, a strict-inequality version cannot be substituted without checking its scope and the endpoint \(r=9\).
\subsection{Short English statement}
For each \(r\geq 9\), the question asks whether the Nagata conjecture alone for the plane blown up at \(r\) general points guarantees an ample line bundle whose Seshadri constant is irrational at a very general point. The ground field, the meaning of general centers, and the precise Nagata hypothesis remain unspecified. The supplied record labels the question open; this draft does not independently verify that status.
\subsection{Sources}
\begin{itemize}
\item Local project registry: supplied UnsolvedMath source record, permanent ID \(30003978\), source problem code \texttt{catalogue entry}, title ``Irrational Seshadri Constants from the Nagata Conjecture.'' The record belongs to the Oberwolfach Reports Open Problems collection. No individual proposer is supplied.
\item Supplied primary report citation: \emph{Mini-Workshop: Asymptotic Invariants of Homogeneous Ideals} (2019), \emph{Oberwolfach Reports}, pp.\ 2739--2767, DOI \texttt{10.4171/owr/2018/45}. \url{https://doi.org/10.4171/owr/2018/45}. The record separately gives proposed year 2018 and the tag \texttt{owr-year:2018}; these supplied dates have not been reconciled.
\item Supplied literature source: ``Rationality of Seshadri constants.'' \url{https://doi.org/10.1016/j.jpaa.2020.106345}.
\item Supplied literature source: ``Irrational Seshadri constants.'' \url{https://www.cmi.ac.in/~krishna/irrational-seshadri-constants.pdf}.
\item Supplied literature assessment dated exactly \texttt{2026-08-21}: the record labels the question open and reports that later work provides conditional results rather than a proof from Nagata alone. This is supplied dated curation derived from the OWR material, not an independently verified assessment. The supplied background explicitly requests independent verification of time-sensitive status and cited claims.
\item The registry assigns statement status \texttt{corrected\_verified}, explaining that notation, typography, and sentence structure were repaired while preserving mathematical content. This is supplied editorial metadata, not independent mathematical verification by this draft.
\end{itemize}
\end{document}
